\documentclass{article}

\usepackage{../handouts}

\title{CSCI 432 Handout 12: Graph Algorithm Recap}
\author{Name: \rule{3in}{0.15mm}}

\begin{document}
\maketitle

\section*{Shortest Paths}

Input: A (directed or undirected)
weighted graph such that all edge weights are positive,
$G=(V,E,\omega)$, where $\omega \colon E \to R_{+}$, a start node $s$, and an
end node $t$.

Output: A shortest path between $s$ and $t$.

\begin{enumerate}
    \item Give an example of a graph such that the shortest path IS NOT unique..

        \textbf{Answer}
        \vspace{1in}

    \item Give an example of a graph such that the shortest path IS unique..

        \textbf{Answer}
        \vspace{1in}

    \item Give an example of a graph for which there is no path between $s$ and
        $t$.

        \textbf{Answer}
        \vspace{1in}

\end{enumerate}

\newpage
\section*{Dijkstra's Algorithm}

Let's remember how Dijkstra's algorithm works.

Initial set-up: We start create a priority queue where each vertex has priority
$\infty$, except $s$, which has priority $0$.

The main loop: We pop off the next element $x$ of the priority queue, as we now
know it's distance to the start $s$. We then perform the `relaxation' step of
updating all priorities to allow for going through $x$. If going through $x$ is
shorter than tentative distance ($priority(x)+w(x,v)$ < current priority), then
we update the priority.

When $t$ is popped, we return the priority of $t$.

\begin{enumerate}
    \item Walk through this algorithm for finding the shortest distance
        between $0$ and $6$ in the following graph:
        \includegraphics[height=3in]{graph}

    \item What is the runtime of Dijsktra's algorithm?
        \textbf{Answer}
        \vspace{2in}

    \item Write psuedo-code so that the algorithm returns a shortest path
        between $s$ and $t$, not just the distance.
        \textbf{Answer}
        \vspace{2in}

    \item True or false: If we don't stop when we pop $t$, then we wind up with
        a minimum spanning tree.
        \textbf{Answer}
        \vspace{2in}

    \item What is the loop invariant for Dijkstra's algorithm
        (here, just think of the loop for the
        algorithm that returns the distance)?
        \textbf{Answer}
        \vspace{2in}
\end{enumerate}


\newpage
\section*{Other Well-Known Graph Algorithms}

The following algorithms are well-known graph algorithms.
For each algorithm, (a) recall or lookup the algorithm, maybe writing pseudocode
and identifying the runtime; (b) work through a small
example; (c) identify the loop invariant.

\begin{enumerate}
    \item Bellman--Ford: This algorithm computes the shortest path from $s$ to
        all other vertices. Unlike Dijkstra's algorithm, it allows negative edge
        weights (and can detect negative cycles).
        The initialization of Bellman--Ford
        is similar to that of Dijkstra (all
        vertices are labeled with $\infty$ as the estimated distances, except
        $s$ is labeled with $0$). Bellman--Ford does not use a priority queue
        though, it relaxes every edge every time ... stopping after $|V|$
        iterations.
    \item Floyd--Warshall algorithm: Another algorithm for computing shortest
        paths, allows negative weights (but no negative cycles).  This algorithm
        uses dynamic programming to compute all pairs shortest paths.  Along the
        way, this algorithm computes, for each pair $(v,w)\in V^2$ and each $k
        \in \{1, 2, \ldots, |V|\}$, the length of the shortest path from $v$ to
        $w$ only using the first $k$ vertices.
    \item Prim's Algorithm. Computes the MST of a weighted (undirected) graph.
        Similar to Dijkstra, it starts from an arbitrary node and uses a
        priority queue.
\end{enumerate}

\end{document}
